\subsection{The graph terms of the doubled marked family}
\label{dg:section}

We compute the graph functional in the rational symplectic basis in
which the undoubled sections are defined. Write $g=r|\Psi|$,
$\Omega_x$ for the Siegel period matrix, $Y_x=\operatorname{Im}\Omega_x$,
and $h_x$ for the inverse of the first period matrix. Thus a real
period vector with symplectic coordinates $(u,v)$ has normalized
complex period $u+\Omega_x v$. Use Lebesgue measure on the $ng$ real
position coordinates and mass one for their integral finite lattice.
For one row the Gaussian, apart from the determinant, is
\[
 q_{x,z}(s)=\exp\!\left(\pi i\,s^t\Omega_xs+
          2\pi i\,s^t h_xz+\tfrac{\pi}{2}
                   z^t h_x^tY_x^{-1}h_xz\right).
\]
Put $Q_{x,z}=(\det h_x)^{n/2}\prod_{k=1}^nq_{x,z_k}(s_k)$
and $Q_{x,P}=P(\partial_z)Q_{x,z}|_{z=0}$. These give exactly
the rational theta sums defining $\theta_P(\phi;x)$ with finite
test $\phi$. We use inner products linear in their first variable.

\begin{lemma}\label{dg:gaussian}
For the family and weight of Proposition~\ref{arrep:reproduce},
\[
 \sum_j\langle Q_{x,P_j},Q_{y,P_j}\rangle_\infty
 =c_\infty K(x,y),\qquad
 c_\infty=2^{-ng/2}\prod_{\tau\in\Psi}
     \bigl(|\det G_\tau|\,[-\operatorname{Im}\tau(\kappa_0)]^r
                                                        \bigr)^{n/2}.
\]
This factor is independent of the weight, the polynomials, and both
frames.
\end{lemma}
\begin{proof}
First we show that $Q_{x,P}$ is holomorphic in $x$ with values in
Schwartz space. The physical period map at $\tau$ has components
$\langle x_i,s\rangle$ for $i\le p$ and $(x^{-1}s)_i$ for $i>p$.
Thus $h_x,\Omega_x$ are holomorphic. For the remaining quadratic put
$C=h_x^tY_x^{-1}h_x$. Choose a purely imaginary $\theta\in E$
which generates $E/\mathbf Q$. Such elements exist: writing the
purely imaginary subspace as $\eta E^+$ for any nonzero imaginary
$\eta$, distinct embeddings have distinct restrictions to this
subspace (evaluate first on $\eta$, then on its multiples).
Finitely many proper rational linear subspaces therefore suffice to
exclude all coincidences. Multiplication by $\theta$ in physical
period coordinates is diagonal $D$, with entries $\bar\tau(\theta)$
on the positive rows and $\tau(\theta)$ on the negative rows.
Write its rational blocks in the symplectic basis as
$\left(\begin{smallmatrix}a&b\\ c&d\end{smallmatrix}\right)$.
Applying it to first periods gives $h_xD=(a+\Omega_xc)h_x$.
Multiplication by $\theta$ is skew-adjoint for the alternating trace
pairing and complex-linear on the normalized periods. Hence it is
skew-adjoint for their Hermitian form with matrix $Y_x^{-1}$: the
imaginary part is the alternating pairing, and replacing one argument
by $i$ times that argument recovers the real part. With $A=a+\Omega_xc$,
this gives $A^tY_x^{-1}+Y_x^{-1}\bar A=0$. Consequently
\[
                  D^t C+C D=2i h_x^t c h_x .
\]
Whenever $D_i+D_j\ne0$, $C_{ij}$ is holomorphic by division.
The remaining pairs are exactly the opposite signs at the same
place. Their terms in the product exponential multiply the other
terms by exponentials of mixed row contractions. They disappear
under $P(\partial_z)$ by the mixed Laplacian, as in
Proposition~\ref{arrep:reproduce}. On each compact subset of
$\mathcal D$, $Y_x$ is bounded below by a positive matrix, so the
resulting polynomial Gaussians and their derivatives have locally
uniform Schwartz bounds. This proves holomorphy and hence
sesquiholomorphy of the displayed pairing.

On replacing $x$ by $xb$, with triangular diagonal blocks $A_0,D_0$,
the physical period map acquires left factor $A_0^t\oplus D_0^{-1}$.
Thus $\Omega_x$ is unchanged, and $h_x$ gains right factor
$A_0^{-t}\oplus D_0$. Ordinary differentiation and the row-highest
law give $Q_{xb,P}=\chi_\lambda(b)Q_{x,P}$. It suffices to compare
the two kernels on the diagonal of orthoframes and then on the full
diagonal by triangular equivariance.

For an orthoframe $x$, put $d_\tau=-\operatorname{Im}\tau(\kappa_0)$.
In its basis the physical periods of a vector $s$ are the conjugates
of its positive coefficients and its negative coefficients themselves.
The imaginary part of the Hermitian pairing with diagonal entries
$d_\tau$ in these coordinates is therefore the original alternating
form $\operatorname{Tr}_{E/\mathbf Q}(-\kappa_0\langle s,s'\rangle/2)$.
In normalized period coordinates the same form is
$\operatorname{Im}((u+\Omega_x v)^tY_x^{-1}
\overline{(u'+\Omega_x v')})=v^tu'-u^tv'$. Equality of imaginary
parts gives
$h_x^tY_x^{-1}\bar h_x=\operatorname{diag}(d_\tau 1_r)_\tau$.
Completing the real Gaussian square now gives
\[
 \langle Q_{x,z},Q_{x,\bar w}\rangle_\infty
 = |\det h_x|^n\det(2Y_x)^{-n/2}
          \exp\!\left(\pi\sum_{\tau,i,k} d_\tau z_{\tau ik}w_{\tau ik}\right).
\]
This is the orthoframe Gaussian of Proposition~\ref{arrep:reproduce}
times $c_\infty$, since $|\det x_\tau|^2=|\det G_\tau|^{-1}$.
Differentiate and use triangular equivariance. Sesquiholomorphic
kernels equal on the full diagonal are equal everywhere: the power
series of their difference at each diagonal point has coefficients
detected by independent Wirtinger derivatives on the diagonal,
and $\mathcal D$ is connected.
\end{proof}

\begin{proposition}\label{dg:identity}
Use the marked family in Lemmas~\ref{bfam:genus} and~\ref{bfam:phase},
even $n>4r$, and the column-invariant polynomial family of
Proposition~\ref{arrep:reproduce}. Let $\sum_\Lambda w_\Lambda$
denote the finite genus sum with reciprocal stabilizer weights.
Write $\phi_{L,\Lambda}$ for the marked vector on row $L$,
$\Delta$ for its prescribed rational row group, and
$w(\gamma)\phi_{L,\Lambda}=\mu(\gamma,L)\phi_{\gamma L,\Lambda}$.
For all permissible rows $L_1,L_2$ and all $x,y\in\mathcal D$,
\[
\begin{split}
 &\sum_\Lambda w_\Lambda\sum_j
    \theta_{P_j}(\phi_{L_1,\Lambda};x)
           \overline{\theta_{P_j}(\phi_{L_2,\Lambda};y)}\\
 &\quad=c_\infty\sum_{\gamma\in\Delta}\overline{\mu(\gamma,L_2)}
               K(x,\gamma y)
       \sum_\Lambda w_\Lambda
       \langle\phi_{L_1,\Lambda},\phi_{\gamma L_2,\Lambda}\rangle_f .
\end{split}
\]
The graph sum is absolutely convergent. This is the kernel identity
for this genus, including the phase of each marked row.
\end{proposition}
\begin{proof}
In the doubled oscillator use the rational diagonal polarization.
For a pure tensor $u\otimes\bar v$ the change to this polarization
takes difference and midpoint of the two position coordinates,
followed by Fourier transform in the midpoint. Evaluation at the
origin is therefore $\langle u,v\rangle$ in the self-dual position
measure, both at infinity and at the finite places. In the source
orientation the coordinate function is $W_{u,v}(-s)$ as in
Lemma~\ref{bfam:phase}. The product of the two rational theta sums
is the rational theta sum after this change: the rational pairs of
positions are in bijection with their rational differences and
midpoints, and Poisson summation in the latter variable replaces
it by its rational dual. This proof works for arbitrary Schwartz
inputs. In particular there is no primitivity restriction on the
rational differences. A common rescaling of the finite intertwiner
and the inverse rescaling of the archimedean intertwiner leave the
global evaluations unchanged.

Here it matters that the inserted archimedean test
$\sum_j Q_{x,P_j}\otimes\overline{Q_{y,P_j}}$ is invariant under the
compact column group in the doubled model. In physical period
coordinates a column unitary $u_\tau$ acts by $u_\tau$ in negative
rows and by $\bar u_\tau$ in positive rows. It preserves the
polarization, so its oscillator action fixes the Gaussian vacuum
line at each frame. One can see this by the first-order Gaussian
equations $(\partial_s-2\pi i\Omega_xs)q_{x,0}=0$: the operators
are the positive-polarization annihilators and their common solution
line is spanned by this Gaussian. Its unit multiplier is the same
at any two frames. Indeed their vacuum pairing is nonzero by the
Gaussian integral with matrix $-i(\Omega_x-\bar\Omega_y)$ of
positive real part. To check the action on the inserted vectors,
the Weyl translate of the vacuum with shift $v$ and modulation $u$
is $q_{x,z}\exp(-\pi\zeta^tY_x^{-1}\bar\zeta/2)$, where
$\zeta=u+\Omega_xv=h_xz$. This follows by replacing $s$ with $s+v$
and multiplying by $\exp(2\pi i u^t(s+v/2))$. The column action
preserves the Hermitian norm in this formula. Thus it transforms
the coherent parameters by the stated linear substitutions and the
common unit multiplier. The multipliers cancel between the two
copies, and column invariance of the sum of differential symbols
proves the assertion (the transposes in differentiation are again
paired unitary substitutions).

Apply Lemma~\ref{bfam:genus} to this invariant Schwartz test.
At the Siegel--Weil point the section for a Schwartz function $\Phi$
is $g\mapsto(\omega(g)\Phi)(0)$. Here the doubled group has split
rank $r$, the positive definite column space has dimension $n$ and
the parameter is $s_0=(n-r)/2$. The theta integral uses Haar
probability, giving coefficient one. Multiplying back by the
total mass gives the same weights on both sides.
The splitting character on the Siegel Levi can be taken trivial:
its restriction to $\mathbb A_{E^+}^\times$ is required to be the
$n$-th power of the quadratic extension character, and $n$ is even.
In this splitting the diagonal column group acts by substitution.
The same reference gives the doubled tensor intertwiner and its
value at zero; its general second-factor twist by the splitting
character is trivial here.
\cite[\S\S2.8, 3.1, 11.3 and Theorem~7.1]{ganQiuTakeda}

We specify the graph convention. Label the graph by the rational
row isometry $\gamma$ acting on the second member of the tensor to
give the diagonal functional. Thus for the seed graph one evaluates
at zero, and for this graph one first translates that member by
$\gamma$. In row coordinates a representative is the conjugate of
$\operatorname{diag}(1,\gamma^{-t})$ by the rational change of
polarization with blocks
$\left(\begin{smallmatrix}1&\kappa_0G/4\\-1&\kappa_0G/4\end{smallmatrix}\right)$.
Equivalently the action on the second undoubled theta test is the
rational oscillator action carrying its frame $y$ to $\gamma y$.
To check the scalar, let
$\left(\begin{smallmatrix}a&b\\c&d\end{smallmatrix}\right)$ be
the induced Siegel transformation in the symplectic period basis,
$U=c\Omega_y+d$, $\Omega'=(a\Omega_y+b)U^{-1}$.
The equality of the period maps on $\gamma s$ at $\gamma y$ and on
$s$ at $y$ gives $h_{\gamma y}=U^{-t}h_y$,
$\Omega_{\gamma y}=\Omega'$. The Gaussian action sends
$Q_{y,z}$ to $Q_{\gamma y,z}$ up to a constant phase independent
of $y,z$. Indeed under a real linear coordinate change this follows
by substitution, and under adding a real symmetric matrix to the
period it follows by multiplying by the corresponding quadratic
phase. Under Fourier inversion use the Gaussian integral
\[
 \int_{\mathbb R^g}e^{\pi i s^t\Omega s+2\pi i s^t t}\,ds
       =\det(-i\Omega)^{-1/2}e^{-\pi i t^t\Omega^{-1}t}.
\]
For the Fourier generator $\Omega'=-\Omega^{-1}$, the factor
$e^{-\pi i\zeta^t\Omega^{-1}\zeta}$ combines with the original
quadratic correction to give the one at
$\zeta'=\Omega^{-t}\zeta$, since
$\zeta'^tY'^{-1}\zeta'=\zeta^tY^{-1}\zeta-
2i\zeta^t\Omega^{-1}\zeta$. The determinant in $Q$ absorbs the
period-dependent determinant after taking $n$ rows. These generators
generate the rational symplectic group by symplectic elimination.
Choose the reciprocal scalar on the finite operator. Its action
$w(\gamma)$ is then exactly the one characterized by
$\theta(w(\gamma)\phi;\gamma y,z)=\theta(\phi;y,z)$: for the
Fourier generator this is rational Poisson summation, for rational
linear changes it is reindexing with the product formula, and for
the shear it is cancellation of the rational quadratic phases.
This characterization determines the finite operator, since at
fixed $y$ the distinct rational position frequencies of the theta
series can be separated by Fourier coefficients in $h_yz$ after
removing the common nonzero Gaussian prefactor.

These are also the normalizations of the graph section. Indeed the
diagonal-polarization change intertwines the rational row actions
with their conjugates in the doubled group. For the linear changes,
quadratic phases and Fourier generators this is just covariance
of Weyl translations under these operations. Any scalar ambiguity
between the two intertwiners is independent of the test by
irreducibility of the Schrödinger representation; its global value
is one since both rational actions preserve the nonzero rational
theta distribution by Poisson summation. The choice of another
rational Siegel-parabolic representative likewise preserves the
global evaluation, by the product formula on its determinant and
triviality of its unipotent character at zero. The term on the
graph consequently is precisely
\[
 \langle\phi_{L_1,\Lambda},w(\gamma)\phi_{L_2,\Lambda}\rangle_f\,
          \sum_j\langle Q_{x,P_j},Q_{\gamma y,P_j}\rangle_\infty .
\]
Graphs outside $\Delta$ vanish by the marked support condition.
For $\gamma\in\Delta$ the inner product in the second variable
contributes $\overline{\mu(\gamma,L_2)}$.

For a nongraph there is a nonzero rational isotropic row line
$E\xi$ in the intersection of its maximal isotropic with the first
copy. Indeed either projection being singular forces the other to
be singular by orthogonality and dimension. For every $v\in E^n$
the finite Weyl translation on the first input by $v\otimes\xi$ acts,
after passage to this polarization, by a pure modulation. It takes
value one at the origin: in difference and dual coordinates the
action is $\phi(t-d)\psi_f(\langle b,t-d/2\rangle)$ with $d=0$.
There is no residual scalar, since linear symplectic covariance
preserves the symmetric Weyl operators including their half
pairings. The resulting finite evaluation $\ell$ thus satisfies
$\ell(\rho(v\otimes\xi)\phi_1,\phi_2)=\ell(\phi_1,\phi_2)$.
The marked isotropic relation has finite coefficients $a_v$ with
$\sum a_v\ne0$ and $\sum a_v\rho(v\otimes\xi)\phi_1=0$.
It gives $(\sum a_v)\ell(\phi_1,\phi_2)=0$ before any
archimedean evaluation. This accounts for every nongraph.

Finally the polynomial differentiations are legitimate in this
expansion. On compact sets of complex $z,w$ the archimedean
Gaussians, with every fixed parameter derivative, are bounded in
Schwartz space, by completion of the square and $Y_x,Y_y>0$.
The same holds after the fixed partial Fourier transform.
For a Siegel representative with Iwasawa decomposition $pk$ and
Levi block $a_p$, evaluation is bounded by
$|\det a_p|_{\mathbb A_E}^{n/2}
\sup_k|(\omega(k)\Phi)(0)|$; the supremum is uniform on these
bounded test families by continuity and compactness. The finite
tests lie in finitely many compact orbits under the maximal compact
subgroup (each smooth test has open stabilizer there).
The scalar Siegel Eisenstein majorant converges absolutely for
$\operatorname{Re}s>r/2$, which contains $s_0$ since $n>4r$.
Thus either Cauchy's formula on a slightly larger parameter
polydisc or direct Schwartz differentiation permits the fixed
finite insertions under the sum. The genus formula above was only
needed for the invariant sum of inserted vectors. Gaussian theta
series on the left also converge normally since the finite test
restricts rational positions to a fixed fractional lattice.
Combining the surviving graph evaluations with Lemma~\ref{dg:gaussian}
proves the identity.
\end{proof}

\src{src/Development/AbelianVarietyTheory/Proofs/SiegelThetaTransformation/EisensteinSeriesTermUnitaryGraph.lean}{33}{The graph representative and both translated tests};
\src{src/Development/AbelianVarietyTheory/Proofs/SiegelThetaTransformation/KernelIdentity.lean}{181}{The class-weighted differentiated graph identity}.
